\section{Einleitung}

\subsection{Motivation}
Die Sicherheit von Linux-Systemen ist ein fundamentales Anliegen in der heutigen IT-Landschaft. 
Das Linux Security Module Framework stellt Schnittstellen bereit, über die Sicherheitsmodule unterschiedliche Zugriffskontrollmodelle im Linux-Kernel durchsetzen können\citep{kernelorgLSM}.
LSM-Hooks ermöglichen es, sicherheitsrelevante Kerneloperationen abzufangen und auf dieser Grundlage Zugriffsentscheidungen zu treffen.
Ein LSM-Hook prüft einen Zugriff jedoch grundsätzlich zu dem Zeitpunkt, an dem die jeweilige Kerneloperation ausgeführt wird.
Bei dauerhaft bestehenden Zugriffen wird eine spätere Änderung der Entscheidungsgrundlage dadurch nicht automatisch erneut bewertet.



\subsection{Wissenschaftliche Fragestellung}
Diese Arbeit untersucht, wie zentrale Prinzipien von Attribute Stream-Based Access Control für Dateiöffnungen unter Linux prototypisch umgesetzt werden können.
Im Mittelpunkt steht dabei die Frage, wie dynamische Attribute aus dem Userspace für Zugriffsentscheidungen eines eBPF-LSM-Programms bereitgestellt und Änderungen bei bereits bestehenden Dateizugriffen nachträglich berücksichtigt werden können.

Daraus ergibt sich folgende wissenschaftliche Fragestellung:

\begin{quote}
Wie kann Attribute Stream-Based Access Control für Dateiöffnungen prototypisch mittels eBPF-LSM unter Linux umgesetzt werden, und welche technischen Einschränkungen ergeben sich insbesondere bei der nachträglichen Kontrolle bereits bestehender Dateizugriffe?
\end{quote}

Zur Beantwortung dieser Fragestellung werden insbesondere die Überführung textueller Policy-Dateien in eine kernelgeeignete Repräsentation, die Bereitstellung dynamischer Attribute für Zugriffsentscheidungen im Kernel sowie die nachträgliche Überwachung bestehender Zugriffssituationen betrachtet.

\subsection{Zielsetzung}
Zielsetzung der Arbeit ist ein Prototyp, der zentrale Eigenschaften von Attribute Stream-Based Access Control auf die Kontrolle von Dateiöffnungen unter Linux überträgt.
Insbesondere sollen Stream-Attribute wie Zeitinformationen sowie externe, zur Laufzeit aktualisierte Attribute in die Zugriffsentscheidung einbezogen werden.
Der Prototyp soll textuelle Policies einlesen, in eine kernelgeeignete Repräsentation übersetzen und neue Dateiöffnungen mittels eBPF-LSM im Kernel bewerten.
Nach der erfolgreichen Aktivierung geänderter Policies oder Attribute sowie beim Erreichen entscheidungsrelevanter Zeitgrenzen soll er bereits bestehende Dateizugriffe im Userspace ereignisbasiert erneut bewerten und bei einer festgestellten Verletzung einen Entzug versuchen.

Nicht Gegenstand der Arbeit ist die vollständige Implementierung der Publish-Subscribe-Architektur der Streaming Attribute Policy Language oder eines allgemeinen, eigenständigen ASBAC-PDP.
Der Prototyp ist auf Dateiöffnungen über den LSM-Hook \texttt{file\_open} und auf die nachträgliche Untersuchung geöffneter File Descriptors beschränkt.


\subsection{Aufbau der Arbeit}
Kapitel~2 führt die für die Arbeit erforderlichen Grundlagen ein. Dazu gehören der Reference-Monitor-Ansatz, Attribute-Based Access Control und Attribute Stream-Based Access Control sowie die für die Umsetzung relevanten Linux-Technologien LSM, eBPF und eBPF-Maps. Zudem wird der eigene Ansatz gegenüber der Publish-Subscribe-Architektur von SAPL und verwandten Konzepten fortdauernder Autorisierung abgegrenzt.

Kapitel~3 leitet aus der Zielsetzung die funktionalen, operativen und entwicklungsbezogenen Anforderungen an den Prototyp ab. Kapitel~4 entwickelt darauf aufbauend die Systemarchitektur und begründet die wesentlichen Entwurfsentscheidungen. Im Mittelpunkt stehen die Aufgabenteilung zwischen Userspace und Kernelspace, die Repräsentation und konsistente Aktualisierung von Policies und Attributen sowie die ergänzende Nachbewertung bereits geöffneter File Descriptors.

Kapitel~5 beschreibt die konkrete Umsetzung des Prototyps. Es erläutert die Struktur des Rust-Workspace, die gemeinsame Map-ABI, die Policy- und Attributloader, den Kernelspace-PEP auf Grundlage des LSM-Hooks \texttt{file\_open}, den Userspace-PEP und das Administrationstool. Darüber hinaus werden Fehlerbehandlung, Build, Deployment und implementierungsbedingte Grenzen dargestellt.

Kapitel~6 evaluiert den Prototyp auf einem dedizierten Linux-Zielsystem. Unit-, Komponenten- und End-to-End-Tests untersuchen das funktionale Verhalten und ausgewählte Fehlerfälle; ergänzende Messungen charakterisieren den Aufwand kontrollierter Dateiöffnungen sowie die Reaktions- und Entzugslatenz. Kapitel~7 beantwortet abschließend die Forschungsfrage, ordnet die Aussagekraft der Ergebnisse ein und zeigt Ansatzpunkte für weiterführende Arbeiten auf.

\newpage
